\chapter{理論的基礎}
\label{c-continuous}

本章では optimal transport の基礎を述べる。まず確率ベクトル $(\a,\b)$ 間の最適マッチングとカップリングの関連概念を導入し、この計算を離散測度 $(\al,\be)$ 間の輸送へ段階的に一般化し、最後に任意の測度の一般的な設定を扱う。初読時には、これらの最後の微妙な点を省略し、確率ベクトル、すなわちヒストグラム間の計算にのみ注目してもよい。これが第~\ref{c-algo-basics}章および第~\ref{c-entropic}章で詳述するアルゴリズムを実装するための唯一の前提条件である。より経験豊富な読者は、任意の測度に適用される定式化を考慮することでこの問題のより深い理解に到達でき、より高度な問題（\eg 点群の位置を移動させるため、あるいは点が連続密度からサンプリングされる統計的な設定において）に適用することが可能になる。

%%%%%%%%%%%%%%%%%%%%%%%%%%%%%%%%%%%%%%%%%%%%
\section{ヒストグラムと測度}

確率単体
\eq{
	\simplex_n \eqdef \enscond{\a \in \RR_+^n}{ \sum_{i=1}^n \a_i = 1 }
}
に属する任意の要素 $\a \in \simplex_n$ に対して、ヒストグラムと確率ベクトルという用語を互換的に用いる。本レビューの大部分は、optimal transport が単体上に誘導する幾何学の研究に専ら焦点を当てる。

%%%%%%%%%%%%%%
\begin{rem1}{離散測度}
重み $\a$ と位置 $x_1,\dots,x_n\in\X$ を持つ離散測度は
\eql{\label{eq-discr-meas}
	\al = \sum_{i=1}^n \a_i \de_{x_i}
}
と表される。ここで $\de_x$ は位置 $x$ における Dirac 測度であり、直観的には位置 $x$ に無限に集中した単位質量である。このような測度は、$\a\in\simplex_n$ のとき確率測度を表し、ベクトル $\a$ のすべての要素が非負のとき、より一般に正測度を表す。質量を持たない位置が考慮される退化した状況を避けるため、離散測度を考える際には $\a$ のすべての要素が正であると仮定する。
\end{rem1}
%%%%%%%%%%%%%%



%%%%%%%%%%%%%%
\begin{rem2}{一般の測度}
OT の便利な特徴は、離散と連続のいずれか、または両方である測度を同一の枠組みの中で扱えることである。そのために、空間 $\X$ 上の Radon 測度の集合 $\Mm(\X)$ を用いる。この集合の形式的な定義には $\X$ が距離（通常 $\dist$ と表記される）を備えていることが必要である。これは測度に連続関数（$f \in \Cc(\X)$ と表記される）に対する積分（「テスト」）を通じてのみアクセスできるためである。

$f \in \Cc(\X)$ の離散測度 $\al$ に対する積分は和として計算される：
\eq{
	\int_\X f(x) \d\al(x) = \sum_{i=1}^n \a_i f(x_i).
}
%
より一般的な測度、例えば $\X=\RR^\dim$（ここで $\dim \in \NN^*$ は次元）上の測度は、Lebesgue 測度に関する密度 $\d\al(x)=\density{\al}(x)\d x$ を持ちうる。しばしば $\density{\al} = \frac{\d\al}{\d x}$ と表記され、これは
\eq{
	\foralls h \in \Cc(\RR^\dim), \quad
	\int_{\RR^\dim} h(x) \d\al(x) =  \int_{\RR^\dim} h(x) \density{\al}(x) \d x
}
を意味する。
%
任意の測度 $\al \in \Mm(\X)$（密度を持つ必要も Dirac の和である必要もない）は、任意の連続関数 $f \in \Cc(\X)$ に対して積分でき $\int_\X f(x) \d\al(x) \in \RR$ を得られるという事実によって定義される。$\X$ がコンパクトでない場合、$f$ がコンパクト台を持つか、少なくとも無限遠で極限が $0$ であることも課す必要がある。
%
測度は、ある意味で関数より「正則性が低い」が、超関数（滑らかな関数の双対として定義される）よりは正則である。例えば、Dirac の導関数は測度ではない。
%
$\X$ 上のすべての正測度の集合を $\Mm_+(\X)$ と表記する。確率測度の集合は $\Mm_+^1(\X)$ と表記され、任意の $\al \in \Mm_+^1(\X)$ が正であり、かつ $\al(\X)=\int_\Xx \d\al = 1$ であることを意味する。
%
図~\ref{fig-measures}は、ヒストグラムを超えて、本書で考慮される異なるクラスの測度の視覚化を提供する。
\end{rem2}
%%%%%%%%%%%%%%


\begin{figure}[h!]
\centering
\begin{tabular}{@{}c@{\hspace{1mm}}c@{\hspace{1mm}}c@{\hspace{1mm}}c@{}}
\includegraphics[width=.24\linewidth]{measures/discrete-1d.pdf}&
\includegraphics[width=.15\linewidth]{measures/discrete-2d}&
\includegraphics[width=.24\linewidth]{measures/density-1d}&
\includegraphics[width=.22\linewidth]{measures/density-2d}\\
離散 $d=1$ & 離散 $d=2$ & 密度 $d=1$ & 密度 $d=2$
\end{tabular}
\caption{\label{fig-measures}
離散分布 $\al = \sum_{i=1}^n \a_i \de_{x_i}$（赤は経験的一様分布 $\a_i=1/n$、青は任意の分布に対応する）と密度 $\d\al(x)=\density{\al}(x)\d x$（紫色）の模式図。1次元と2次元の両方を示す。1次元の離散分布はステムプロット（長さは $\a_i$ に等しい）で表示され、2次元では点群を使用して表示される（その場合、半径が $\a_i$ に等しいか、視覚的により正確な表現として面積が $\a_i$ に等しい）。
}
\end{figure}





%%%%%%%%%%%%%%%%%%%%%%%%%%%%%%%%%%%%%%%%%%%%
\section{割当てと Monge 問題}

コスト行列 $(\C_{i,j})_{i \in \range{n}, j \in \range{m}}$ が与えられ、$n=m$ を仮定するとき、最適割当て問題は $n$ 個の要素の置換の集合 $\Perm(n)$ における全単射 $\si$ であって
\eql{\label{eq-optimal-assignment}
	\umin{\si \in \Perm(n)} \frac{1}{n}\sum_{i=1}^n \C_{i,\si(i)}
}
を解くものを求める。素朴には、集合 $\Perm(n)$ 内のすべての置換を用いて上のコスト関数を評価できる。しかし、その集合のサイズは $n!$ であり、小さな $n$ に対しても巨大である。例えば、$n$ がわずか 70 のとき、その集合は $10^{100}$ を超える要素を持つ~\citep{Dantzig1983}。したがって、この問題は置換の集合上でコスト関数を最適化する効率的なアルゴリズムが存在する場合にのみ解くことができ、これが\S\ref{s-auction}の主題である。

\begin{rem}[一意性] 最適割当て問題が複数の最適解を持ちうることに注意されたい。例えば、$n=m=2$ であり行列 $\C$ が辺の長さ $1$ の2次元正方形の4つの頂点間の対距離行列であるとする（図~\ref{fig-non-unique-matching}の左のプロットに表されるように）。この場合、2つの割当てのみが存在し、いずれも最適である。
\end{rem}


\begin{figure}[h!]
\centering
\includegraphics[width=.8\linewidth]{non-unique-optimal-matching/non-unique-optimal-matching}
\caption{\label{fig-non-unique-matching}
左：測度 $\alpha$ の青い点と測度 $\beta$ の赤い点は対ごとに等距離である。したがって、マッチング $\sigma=(1,2)$（実線）と $\sigma=(2,1)$（点線）のいずれも最適である。右：Monge 写像は青い測度 $\alpha$ を赤い測度 $\beta$ に対応させることができる。重み $\alpha_i$ は各位置に記された円盤の面積に比例して表示される。ここでの写像は $T(x_1)=T(x_2)=y_2$、$T(x_3)=y_3$ であり、$4\leq i\leq 7$ に対しては $T(x_i)=y_1$ である。
}
\end{figure}


%%%%%%%%%
\begin{rem1}{離散測度間の Monge 問題}\label{rem-monge1}
離散測度
\eql{\label{eq-pair-discr}
	\al = \sum_{i=1}^n \a_i \de_{x_i}
	\qandq
	\be = \sum_{j=1}^m \b_j \de_{y_j}
}
に対して、\citeauthor{Monge1781}の問題~\citeyearpar{Monge1781}は、各点 $x_i$ を単一の点 $y_j$ に対応させ、$\al$ の質量を $\be$ の質量に向けて押し出す写像を求める。すなわち、写像 $T:\{x_1,\dots, x_n\}\rightarrow \{y_1,\dots,y_m\}$ は
\eql{\label{eq-monge-constr}
	\foralls j \in \range{m}, \quad
	\b_j = \sum_{ i : \T(x_i) = y_j } \a_i
}
を満たさなければならない。これを $\T_\sharp \al = \be$ と簡潔に書く。$\b$ のすべての要素が正であるため、この写像は必然的に全射である。
%
この写像は、点 $(x,y) \in \X \times \Y$ に対して定義された関数 $\c(x,y)$ によってパラメータ化された輸送コストを最小化すべきであり、
\eql{\label{eq-monge-discr}
	\umin{\T} \enscond{ \sum_{i} \c(x_i,\T(x_i))  }{ \T_\sharp \al = \be }
}
となる。離散的な点間のこのような写像は、もちろんすべての $x$ と $y$ が異なると仮定すれば、添字 $\si : \range{n} \rightarrow \range{m}$（$j=\si(i)$）を用いて符号化でき、質量保存は
\eq{
	\sum_{i \in \si^{-1}(j)} \a_i = \b_j
}
と書かれる。ここで逆像 $\si^{-1}(j)$ は $j$ の原像集合として理解される。$n=m$ でありすべての重みが一様、すなわち $\a_i=\b_j=1/n$ である特殊な場合には、質量保存制約は $\T$ が全単射であることを意味し、$\T(x_i)=y_{\si(i)}$ となり、Monge 問題は最適マッチング問題~\eqref{eq-optimal-assignment}と同値になる。ここでコスト行列は
\eq{
	\C_{i,j} \eqdef \c(x_i,y_j)
}
である。$n\ne m$ の場合、最適性の議論を別にしても、ある離散測度から別の離散測度への Monge 写像は存在しないことがあることに注意されたい。これは重みベクトルが整合しない場合に起こり、ターゲット測度がソース測度より多くの点を持つ場合（$n<m$）には常にそうなる。例えば、図~\ref{fig-non-unique-matching}の右のプロットは $\al$ と $\be$ の間の（最適な）Monge 写像を示すが、$\be$ から $\al$ への Monge 写像は存在しない。
\end{rem1}
%%%%%%%%%


%%%%%%%
\begin{rem2}{押し出し作用素}\label{rem-push-f}
連続写像 $\T : \X \rightarrow \Y$ に対して、対応する押し出し作用素 $\T_\sharp : \Mm(\X) \rightarrow \Mm(\Y)$ を定義する。
%
離散測度~\eqref{eq-discr-meas}に対して、押し出し操作は測度の台に含まれるすべての点の位置を移動させることに単に対応する：
\eq{
	\T_{\sharp} \al \eqdef \sum_i \a_i \de_{\T(x_i)}.
}
より一般的な測度、例えば密度を持つ測度に対して、押し出しの概念は確率測度の空間的変形（または輸送）を記述する上で基本的な役割を果たす。形式的な定義は以下の通りである。

\begin{defn}[押し出し]\label{defn-pushfwd}
$\T : \X \rightarrow \Y$ に対して、ある $\al \in \Mm(\X)$ の押し出し測度 $\be = \T_\sharp \al \in \Mm(\Y)$ は
\eql{\label{eq-push-fwd}
	\foralls h \in \Cc(\Y), \quad \int_\Y h(y) \d \be(y) = \int_\X h(\T(x)) \d\al(x).
}
を満たす。同値な表現として、任意の可測集合 $B \subset \Y$ に対して
\eql{\label{eq-equiv-pushfwd}
	\be(B) = \al( \enscond{x \in \X}{\T(x) \in B} ) = \al( T^{-1}(B) ).
}
$\T_\sharp$ は正値性と全質量を保存するので、$\al \in \Mm_+^1(\X)$ ならば $\T_\sharp \al \in \Mm_+^1(\Y)$ であることに注意されたい。
\end{defn}

直観的には、可測写像 $T: \X\rightarrow \Y$ はある可測空間から別の可測空間へ単一の点を移動させる関数として解釈できる。$T_\sharp$ は $T$ の拡張であり、$\X$ 上の確率測度全体を $\Y$ 上の新しい確率測度へ移動させることができる。作用素 $T_\sharp$ は、$\X$ 上の測度 $\al$ の各微小質量に写像 $T$ を適用して $\Y$ における微小質量を得ることで、測度を\emph{押し出す}。押し出し作用素 $\T_\sharp : \Mm_+^1(\X) \rightarrow \Mm_+^1(\Y)$ は、$\X$ 上の2つの測度 $\al_1,\al_2$ に対して $T_\sharp(\al_1+\al_2)=T_\sharp\al_1+ T_\sharp\al_2$ という意味で\emph{線形}であることに注意されたい。
\end{rem2}
%%%%%%%

%%%%%%%
\begin{rem2}{多変量密度に対する押し出し}
$\RR^\dim$ 上で密度 $(\density{\al},\density{\be})$ を持つ測度に対して、公式~\eqref{eq-push-fwd}で（$\T$ が滑らかかつ全単射であると仮定して）変数変換を明示的に行うと、押し出しは積分公式における変数変換として密度に線形に作用することがわかる。実際、
\eql{\label{eq-pfwd-density}
	\density{\al}(x) = |\det(\T'(x))|  \density{\be}(\T(x))
}
が成り立つ。ここで $\T'(x) \in \RR^{\dim \times \dim}$ は $T$ の Jacobi 行列（$T$ の各座標の勾配を並べた行列）である。
これより
\eq{
	|\det(\T'(x))| = \frac{ \density{\al}(x) }{ \density{\be}(\T(x)) }
}
が得られる。
\end{rem2}
%%%%%%%


%%%%%%%
\begin{rem2}{任意の測度間の Monge 問題}\label{rem-monge2}
Monge 問題~\eqref{eq-monge-discr}は、2つの任意の確率測度 $(\al,\be)$ が2つの空間 $(\X,\Y)$ 上に台を持ち、写像 $\T : \X \rightarrow \Y$ によって結びつけられる場合に拡張できる。この写像は
\eql{\label{eq-monge-continuous}
	\umin{\T} \enscond{ \int_{\X} \c(x,\T(x)) \d \al(x)  }{  \T_\sharp \al = \be }
}
を最小化する。制約 $\T_\sharp \al = \be$ は、$\T$ が注意~\ref{rem-push-f}で定義された押し出し作用素を用いて $\al$ の質量を $\be$ に押し出すことを意味する。
\end{rem2}
%%%%%%%


%%%%%%%
\begin{rem2}{押し出しと引き戻し}
測度の押し出し $\T_\sharp$ は、関数の引き戻し $\T^\sharp : \Cc(\Y) \rightarrow \Cc(\X)$（関数間の「ワーピング」に対応する）と混同してはならない。引き戻しは、$g \in \Cc(\Y)$ に対して $\T^\sharp g = g \circ \T$ を対応させる線形写像として定義される。押し出しと引き戻しは実際には互いに随伴であり、
\eq{
	\foralls (\al,g) \in \Mm(\X) \times \Cc(\Y), \quad
	\int_\Y g \d( \T_\sharp\al ) = \int_\X (\T^\sharp g) \d\al
}
が成り立つ。$(\al,\be)$ が固定された測度（\emph{例えば} $\RR^\dim$ 上の Lebesgue 測度）に関して密度 $(\density{\al},\density{\be})$ を持つ場合でも、$\T_\sharp \al$ は $\T^\sharp \density{\be}$ を密度として持たないことに注意されたい。これは~\eqref{eq-pfwd-density}における Jacobian の存在による。
%
このことは、OT を画像レジストレーションに用いる際に注意が必要である理由を説明する。OT は画像ワーピング手法として機能するわけではないためである。
%
図~\ref{fig-push-pull}はこれらの押し出しと引き戻し作用素の区別を示している。
\end{rem2}
%%%%%%%


\begin{figure}[h!]
\centering
\begin{tabular}{@{}c@{\hspace{5mm}}c@{}}
\includegraphics[width=.3\linewidth]{push-pull/push-forward}&
\includegraphics[width=.3\linewidth]{push-pull/pull-back}\\
測度の押し出し & 関数の引き戻し
\end{tabular}
\caption{\label{fig-push-pull}
押し出し作用素 $\T_\sharp$（任意の測度を入力として取ることができる）と引き戻し作用素 $\T^\sharp$（関数、特に密度に作用する）の比較。
}
\end{figure}



%%%%%%%%%%%%%%
\begin{rem2}{測度と確率変数}\label{rem-meas-random}
	Radon 測度は確率変数の分布を表すものとしても見ることができる。$\X$ 上の確率変数 $X$ は、実際にはある抽象的な（しばしば明示されない）確率空間 $(\Om,\PP)$ からの写像 $X : \Om \rightarrow \X$ であり、その分布 $\al$ は $\PP(X \in A) = \al(A)=\int_A \d\al(x)$ を満たす Radon 測度 $\al \in \Mm_+^1(\X)$ である。
	%
	同値に、それは $\PP$ の $X$ による押し出し $\al=X_\sharp\PP$ である。
	%
	$\T : \X \rightarrow \Y$ に対する別の押し出し $\be = \T_\sharp\al$ を~\eqref{eq-push-fwd}に従って適用することは、別の確率変数 $Y=\T(X) : \om \in \Om \rightarrow \T(X(\om)) \in Y$ を定義することと同値であり、$\be$ は $Y$ の分布である。
	%
	$Y$ からランダムサンプル $y$ を抽出することは、$X$ から抽出した $x$ を用いて $y=\T(x)$ を計算することで単に達成される。
\end{rem2}
%%%%%%%%%%%%%%


%%%%%%%%%%%%%%%%%%%%%%%%%%%%%%%%%%%%%%%%%%%%
\section{Kantorovich 緩和}\label{sec-kantorovich-relaxation}

割当て問題、およびその注意~\ref{rem-monge1}で提示された Monge 問題における一般化は、実用的な問題に現れるような離散測度の研究に必ずしも適切であるとは限らない。実際、割当て問題は置換問題として定式化されるため、\emph{同じ}サイズの\emph{一様}なヒストグラムの比較にしか使えない。非一様な重みを持つ離散測度への直接的な一般化は Monge の押し出し写像の形式主義を用いて行うことができるが、その定式化は質量保存制約~\eqref{eq-monge-constr}を満たす実行可能解が存在しない場合に退化することもある（注意~\ref{rem-monge1}の末尾を参照）。加えて、割当て問題~\eqref{eq-monge-discr}は組合せ的であり、Monge 問題~\eqref{eq-monge-continuous}の実行可能集合は、質量保存制約を満たすすべての押し出し測度からなる集合として連続的にパラメータ化されているにもかかわらず、\emph{非凸}である。したがって、いずれも元の定式化のままでは解くことが困難である。

\citet{Kantorovich42}の鍵となる着想は、輸送の決定論的な性質、すなわちソース点 $x_i$ が別の点または位置 $y_{\sigma_i}$ または $T(x_i)$ にのみ割り当てられるという事実を緩和することである。Kantorovich は代わりに、任意の点 $x_i$ の質量を複数の位置に分配できることを提案する。Kantorovich は質量輸送が\emph{決定論的}であるべきという考えから離れ、\emph{確率論的}な輸送を考慮する。これはソースから複数のターゲットへの\emph{質量分割}として現在一般に知られているものを許容する。この柔軟性は、置換 $\sigma$ や写像 $T$ の代わりに、カップリング行列 $\P  \in \RR_+^{n \times m}$ を用いて符号化される。ここで $\P_{i,j}$ はビン $i$ からビン $j$ へ、または離散測度の形式主義~\eqref{eq-pair-discr}における $x_i$ から $y_j$ へ流れる質量の量を表す。許容カップリングは Monge 写像よりもはるかに単純な特徴付けを持つ：
\eql{\label{eq-discr-couplings}
	\CouplingsD(\a,\b) \eqdef \enscond{ \P \in \RR_+^{n \times m} }{
		\P \ones_m = \a \qandq
		\transp{\P} \ones_n = \b
	},
}
ここで以下の行列-ベクトル記法を用いた：
\eq{
	\P \ones_m = \left(\sum_j \P_{i,j}\right)_i \in \RR^n
	\qandq
	\transp{\P} \ones_n = \left(\sum_i \P_{i,j}\right)_j \in \RR^m.
}
行列の集合 $\CouplingsD(\a,\b)$ は有界であり $n+m$ 個の等式制約で定義されるため、凸多面体（行列の有限集合の凸包）である~\citep[\S8.1]{brualdi2006combinatorial}。

%
加えて、Monge 定式化（図~\ref{fig-non-unique-matching}の右のプロットに示されるように）が本質的に非対称であったのに対し、Kantorovich の緩和された定式化は常に対称である。すなわち、カップリング $\P$ が $\CouplingsD(\a,\b)$ に属する必要十分条件は $\transp{\P}$ が $\CouplingsD(\b,\a)$ に属することである。Kantorovich の optimal transport 問題は
\eql{\label{eq-mk-discr}
	\MKD_{\C}(\a,\b) \eqdef
	\umin{\P \in \CouplingsD(\a,\b)}
		\dotp{\C}{\P} \eqdef \sum_{i,j} \C_{i,j} \P_{i,j}
}
と定められる。これは線形計画問題である（第~\ref{c-algo-basics}章を参照）。線形計画問題の通常の場合と同様に、最適解は必ずしも一意ではない。

\begin{rem}[鉱山と工場]\label{rem-kantorovich-primal} Kantorovich 問題は、以下の資源配分問題において非常に自然な例示を見出す（\citet{Hitchcock41}も参照）。ある事業者が $n$ 個の倉庫と $m$ 個の工場を運営していると仮定する。各倉庫は工場が適切に稼働するために必要な貴重な原材料を含む。より正確には、各倉庫は整数 $i$ で添字付けされ、$\a_i$ 単位の原材料を含む。これらの原材料はすべて工場に移動させなければならず、工場 $j$ が適切に機能するためには所定量 $\b_j$ が必要である。倉庫 $i$ から工場 $j$ へ資源を移送するために、事業者は輸送会社を利用でき、その会社は位置 $i$ から位置 $j$ へ資源の1単位を移動させるのに $\C_{i,j}$ を請求する。輸送会社は輸送の独占権を持ち、経済のすべての関係者に同じ線形価格体系を適用すると仮定する：資源の $a$ 単位を $i$ から $j$ へ輸送するコストは $a\times \C_{i,j}$ に等しい。

上述の問題に直面して、事業者は式~\eqref{eq-mk-discr}で記述された線形計画問題を解き、各ペア $i,j$ について倉庫 $i$ から工場 $j$ へ輸送すべき物資の量 $\P_{i,j}$ を定量する輸送計画 $\P^\star$ を得ることを選択する。事業者はその計画を実行するために輸送会社に合計 $\dotp{\P^\star}{\C}$ を支払う。
\end{rem}

%%%%%%%
\paragraph{カップリングとしての置換行列}

置換 $\si\in\Perm(n)$ に対して、対応する置換行列を $\P_{\si}$ と書く：
	\eql{\label{eq-perm-matrices}
		\foralls (i,j) \in \range{n}^2, \quad
		(\P_{\si})_{i,j} = \choice{
			1/n \qifq j=\si_i, \\
			0 \quad\text{otherwise.}
		}
	}
この場合に
\eq{\dotp{\C}{\P_{\si}}=\frac{1}{n}\sum_{i=1}^n \C_{i,\si_i}}
が成り立つことが確認でき、割当て問題~\eqref{eq-optimal-assignment}がカップリング $\P$ を正確に置換行列に制限した Kantorovich 問題~\eqref{eq-mk-discr}として書き直せることがわかる：
\eq{\umin{\si \in \Perm(n)} \frac{1}{n}\sum_{i=1}^n \C_{i,\si(i)} = \umin{\si \in \Perm(n)} \dotp{\C}{\P_{\si}}.}
次に、置換行列の集合が \citeauthor{birkhoff} 多面体 $\CouplingsD(\ones_n/n,\ones_n/n)$ に真に含まれることが容易に確認できる。実際、任意の置換 $\si$ に対して $\P_{\si}\ones=\ones_n$ かつ $\transp{\P_{\si}}\ones=\ones_n$ であるが、$\ones_n\transp{\ones_n}/n^2$ は有効なカップリングであるが置換行列ではない。したがって、$\dotp{\C}{\P}$ の最小値は、置換行列のみを考える場合よりもすべての輸送計画を考える場合の方が必然的に小さい：
$$ \MKD_{\C}(\ones_n/n,\ones_n/n) \leq \umin{\si \in \Perm(n)} \dotp{\C}{\P_{\si}}.$$

以下の命題は、これらの問題が実際には同じ最適値を与えること、すなわち2つの一様測度 $\a=\b=\ones_n/n$ の間の Kantorovich 問題~\eqref{eq-mk-discr}を最小化する置換行列が常に存在することを示す。Kantorovich 緩和は割当て問題に対して\emph{タイト}である。
%
図~\ref{fig-matching-kantorovitch}は左側に、この特殊な場合に対応する2次元の最適マッチングの例を示す。
%

\begin{prop}[マッチングに対する Kantorovich の定理]\label{prop-matching-kanto}
	$m=n$ かつ $\a=\b=\ones_n/n$ のとき、問題~\eqref{eq-mk-discr}の最適解 $\P_{\si^\star}$ であって、問題~\eqref{eq-optimal-assignment}の最適置換 $\si^\star \in \Perm(n)$ に対応する置換行列であるものが存在する。
\end{prop}

\begin{proof}
	\citeauthor{birkhoff}の定理~\citeyearpar{birkhoff}は、$\CouplingsD(\ones_n/n,\ones_n/n)$ の端点の集合が置換行列の集合に等しいことを述べる。線形計画法の基本定理~\citep[Theorem 2.7]{bertsimas1997introduction}は、空でない多面体上の線形目的関数の最小値が有限ならば、多面体の端点で達成されることを述べる。
\end{proof}

\begin{figure}[h!]
\centering
\begin{tabular}{@{}c@{\hspace{5mm}}c@{}}
\includegraphics[width=.4\linewidth]{matching-kantorovitch/matching}&
\includegraphics[width=.4\linewidth]{matching-kantorovitch/weighted}
\end{tabular}
\caption{\label{fig-matching-kantorovitch}
%
最適マッチングと一般のカップリングの比較。$x_i$ と $y_j$ の間の黒い線分は、式~\eqref{eq-mk-discr}を解く表示された最適カップリングにおける $\P_{i,j}$ の非零要素を示す。
%
左：最適マッチング。命題~\ref{prop-matching-kanto}の設定（同じ点数 $n=m$ を持つ経験測度）に対応する。
%
右：これら2つの重み付き点群はマッチングできない。代わりに、Kantorovich カップリングを用いて2つの任意の離散測度を対応させることができる。
}
\end{figure}


%%%%%%%
\begin{rem1}{離散測度間の Kantorovich 問題}
形式~\eqref{eq-pair-discr}の離散測度 $\al,\be$ に対して、$\al,\be$ の台の点間のすべての対コストを行列 $\C$ に格納する。すなわち $\C_{i,j} \eqdef \c(x_i,y_j)$ として
\eql{\label{eq-kanto-discr}
	\MK_\c(\al,\be) \eqdef \MKD_{\C}(\a,\b)
}
と定義する。したがって、離散測度間の optimal transport の Kantorovich 定式化は、対応する確率重みベクトル $\a,\b$ の間の問題と同じであるが、コスト行列 $\C$ が $\al$ と $\be$ の台に依存する。ただし、$\MK_\c(\al,\be)$ という記法は、確率重みと台の点の\emph{両方}に関する依存性を明示するため、いくつかの状況で有用である。後者はコスト関数 $\c$ を通じてのみ考慮される。
\end{rem1}
%%%%%%%


\begin{figure}[h!]
\centering
\begin{tabular}{@{}c@{\hspace{5mm}}c@{\hspace{5mm}}c@{}}
\includegraphics[width=.3\linewidth]{settings/discrete}&
\includegraphics[width=.3\linewidth]{settings/semi-discrete}&
\includegraphics[width=.3\linewidth]{settings/continuous}\\
離散 & 半離散 & 連続
\end{tabular}
\caption{\label{fig-settings}
Kantorovich OT の3つの主要なシナリオにおける入力測度 $(\al,\be)$ とカップリング $\Couplings(\al,\be)$ の模式図。第~\ref{c-algo-semidiscr}章が半離散の設定に充てられている。
}
\end{figure}



%%%%%%%
\begin{rem1}{最適割当てとカップリングの応用}
最適輸送計画そのもの（カップリング $\P$ として、あるいは存在する場合は Monge 写像 $\T$ として）は、データサイエンス、特に画像処理において多くの応用を見出している。例えば、コントラスト均等化~\citep{delon2004midway}やテクスチャ合成~\citet{gutierrez2017optimal}に用いられている。
%
OT のイメージングサイエンスへの応用の重要な部分は、画像マッチング~\citep{zhu2007image,wang2013linear,museyko2009application,li2013novel}、画像融合~\citep{courty2016optimal}、医用画像~\citep{wang2011optimal}、形状レジストレーション~\citep{makihara2010earth,lai2014multi,su2015optimal}、画像電子透かし~\citep{mathon2014optimal}に関するものである。天体物理学では、OT は初期宇宙の復元に用いられている~\citep{FrischNaturee}。
%
OT は音楽採譜~\citep{flamary2016optimal}にも用いられており、マッチングデータを解釈するための経済学への多数の応用がある~\citep{galichon2016optimal}。
%
最後に、OT の手法を用いて（またはそれに触発されて）計算された輸送写像は、サンプリング~\citep{reich2013nonparametric,oliver2014minimization}や Bayes 推論~\citep{kim2013efficient,el2012bayesian}の実行にも有用であることを注記しておく。
\end{rem1}
%%%%%%%


%%%%%%%
\begin{rem2}{任意の測度間の Kantorovich 問題} $\MK_\c$ の定義~\eqref{eq-kanto-discr}は、積空間上の同時分布であるカップリング $\pi \in \Mm_+^1(\X \times \Y)$ を考慮することで任意の測度に拡張される。離散的な場合は、この積測度が $\pi = \sum_{i,j} \P_{i,j} \de_{(x_i,y_j)}$ の形であることを課す特殊な状況である。一般的な場合、質量保存制約~\eqref{eq-discr-couplings}は同時確率分布に対する周辺制約として書き直されるべきである：
\eql{\label{eq-coupling-generic}
	\Couplings(\al,\be) \eqdef
	\enscond{
		\pi \in \Mm_+^1(\X \times \Y)
	}{
		P_{\X\sharp} \pi = \al
		\qandq
		P_{\Y\sharp} \pi = \be
	}.
}
ここで $P_{\X\sharp}$ と $P_{\Y\sharp}$ は射影 $P_\X(x,y)=x$ および $P_\Y(x,y)=y$ の押し出し（定義~\ref{defn-pushfwd}を参照）である。
%
図~\ref{fig-settings}は、異なるクラスの問題（離散測度と密度）に対してこれらのカップリング制約がどのように表されるかを示す。
%
\eqref{eq-equiv-pushfwd}を用いると、これらの周辺制約は集合 $A \subset \X$ および $B \subset \Y$ に対して $\pi(A \times \Y)=\al(A)$ および $\pi(\X \times B)=\be(B)$ を課すことと同値である。
%
Kantorovich 問題~\eqref{eq-mk-discr}は
\eql{\label{eq-mk-generic}
	\MK_\c(\al,\be) \eqdef
	\umin{\pi \in \Couplings(\al,\be)}
		\int_{\X \times \Y} \c(x,y) \d\pi(x,y)
}
と一般化される。これは測度の空間上の無限次元線形計画問題である。$(\X,\Y)$ がコンパクト空間で $c$ が連続ならば、常に解が存在することが容易に示せる。実際、$\Couplings(\al,\be)$ は測度の弱位相でコンパクトであり（注意~\ref{dfn-weak-conv}を参照）、$\pi \mapsto \int c \d\pi$ はこの位相に関する連続関数であり、制約集合は空でない（例えば $\al \otimes \be \in \Couplings(\al,\be)$）。
%
図~\ref{fig-couplings}は~\eqref{eq-mk-generic}を解く離散と連続の最適カップリングの例を示す。
%
図~\ref{fig-couplings-simple}は離散と連続の周辺を含む1次元最適カップリングの他の例を示す。
\end{rem2}
%%%%%%%

\begin{figure}[h!]
\centering
\begin{tabular}{@{}c@{\hspace{10mm}}c@{}}
\includegraphics[width=.3\linewidth]{couplings/couplings-continuous}&
\includegraphics[width=.3\linewidth]{couplings/couplings-discr}
\end{tabular}
\caption{\label{fig-couplings}
左：密度を持つ2つの1次元測度間で~\eqref{eq-coupling-generic}を解く「連続」カップリング $\pi$。カップリングは Monge 写像のグラフ $(x,\T(x))$（黒で表示）に沿って局在している。
%
右：形式~\eqref{eq-pair-discr}の2つの離散測度間で~\eqref{eq-mk-discr}を解く「離散」カップリング $\T$。正の成分 $\T_{i,j}$ は位置 $(i,j)$ に $\T_{i,j}$ に比例する半径の黒い円盤として表示される。
}
\end{figure}


\begin{figure}[h!]
\centering
\begin{tabular}{@{}c@{}c@{}c@{}c@{}}
\includegraphics[width=.24\linewidth]{couplings/couplings-1}&
\includegraphics[width=.24\linewidth]{couplings/couplings-2}&
\includegraphics[width=.24\linewidth]{couplings/couplings-3}&
\includegraphics[width=.24\linewidth]{couplings/couplings-4}
\end{tabular}
\caption{\label{fig-couplings-simple}
1次元分布間の最適カップリングの4つの単純な例。上段は写像（矢印）として、下段はカップリングとして表される。\citet{Levy2017review}に触発されている。
}
\end{figure}



%%%%%%%%%%%%%%
\begin{rem2}{確率論的解釈}
Kantorovich 問題は、注意~\ref{rem-meas-random}に従い、確率変数の観点から再解釈できる。実際、問題~\eqref{eq-mk-generic}は
\eql{\label{eq-ot-proba-interpretation}
	\MK_\c(\al,\be) = \umin{(X,Y)} \enscond{ \EE_{(X,Y)}(c(X,Y)) }{ X \sim \al, Y \sim \be }
}
と同値である。ここで $(X,Y)$ は $\X \times \Y$ 上の確率変数の組であり、$X \sim \al$（resp.~$Y \sim \be$）は $X$（resp.~$Y$）の法則が測度として $\al$（resp.~$\be$）でなければならないことを意味する。組 $(X,Y)$ の法則は積空間 $\X \times \Y$ 上の $\pi \in \Couplings(\al,\be)$ である。
\end{rem2}
%%%%%%%%%%%%%%



%%%%%%%%%%%%%%%%%%%%%%%%%%%%%%%%%%%%%%%%%%%%
\section{Optimal Transport の距離的性質}

OT の重要な特徴は、コスト行列が適切な性質を満たせば、ヒストグラムおよび確率測度間の距離を定義することである。実際、OT は点間の ground 距離をヒストグラムまたは測度間の距離に持ち上げる標準的な方法として理解できる。

まず、\citet{RubTomGui00}によって最初に導入された用語を用いて、「ground 距離」行列 $\C$ が固定されている場合を考える。これはビン間の置換コストを表し、比較したい複数のヒストグラムにわたって共有される。以下の命題は、OT がこれらのビン上に台を持つヒストグラム間の有効な距離を提供することを述べる。

\begin{prop}\label{prop-metric-histo}
$n=m$ であり、ある $p \geq 1$ に対して $\C=\distD^p=(\distD_{i,j}^p)_{i,j} \in \RR^{n \times n}$ であると仮定する。ここで $\distD \in \RR_+^{n \times n}$ は $\range{n}$ 上の距離である。\ie
\begin{enumerate}[label=(\roman*)]
	\item $\distD \in \RR_+^{n \times n}$ は対称である；
	\item $\distD_{i,j}=0$ となる必要十分条件は $i=j$ である；
	\item $\foralls (i,j,k) \in \range{n}^3, \distD_{i,k} \leq \distD_{i,j}+\distD_{j,k}$。
\end{enumerate}
このとき
\eql{\label{eq-wass-p-disc}
	\WassD_p(\a,\b) \eqdef \MKD_{\distD^p}(\a,\b)^{1/p}
}
（$\WassD_p$ は $\distD$ に依存することに注意）は $\Si_n$ 上の $p$-Wasserstein 距離を定義する。\ie $\WassD_p$ は対称、正値であり、$\WassD_p(\a,\b)=0$ となる必要十分条件は $\a = \b$ であり、三角不等式
\eq{
	\foralls \a,\b,\VectMode{c} \in \Si_n, \quad \WassD_p(\a,\VectMode{c}) \leq \WassD_p(\a,\b) + \WassD_p(\b,\VectMode{c})
}
を満たす。
\end{prop}

\begin{proof}
距離の対称性と正定値性は容易に証明できる：$\C = \distD^p$ が零の対角を持つため、$\WassD_p(\a,\a)=0$ であり、対応する最適輸送行列は $\P^\star=\diag(\a)$ である。$\distD^p$ のすべての非対角要素の正値性により、$\a\ne \b$ のとき $\WassD_p(\a,\b)>0$ である（この場合、許容カップリングは必ず対角の外側に非零要素を持つため）。$\distD^p$ の対称性により、$\WassD_p(\a,\b)$ 自体が対称関数である。


任意の測度に対する Wasserstein 距離の三角不等式を証明するために、\citet[Theorem 7.3]{Villani03}は所定の構造を持つカップリングの存在を強調する接着補題を用いる。離散的な設定では、この接着カップリングの明示的構成は単純である。
%
$\a,\b,\VectMode{c} \in\simplex_n$ とする。$\P$ と $\Q$ をそれぞれ $\a$ と $\b$ の間、$\b$ と $\VectMode{c}$ の間の輸送問題の2つの最適解とする。
%
$\b$ における零座標から生じうる問題を避けるため、$\b_j>0$ のとき $\tilde\b_j \eqdef \b_j$、そうでないとき $\tilde{\b}_j\eqdef 1$ となるベクトル $\tilde{\b}$ を定義して
\eq{
	\SS \eqdef \P \diag(1/\tilde{\b}) \Q \in \RR_+^{n \times n}
}
とし、$\SS \in \CouplingsD(\a,\VectMode{c})$ であることを確認する：
\eq{
	\SS \ones_n  = \P \diag(1/\tilde{\b}) \Q \ones_n = \P (\b / \tilde{\b}) = \P \ones_{\Supp(\b)} = \a,
}
ここで $\ones_{\Supp(\b)}$ は $\b_j>0$ である添字に1、そうでない添字に0を持つサイズ $n$ のベクトルを表し、$\b_j=0$ である $j$ に対しては必然的に $\P_{i,j} = 0$ であるため $\P \ones_{\Supp(\b)} = \P \ones = \a$ という事実を用いた。同様に $\transp{\SS} \ones_n = \VectMode{c}$ が確認できる。三角不等式は以下から従う：
\begin{align*}
\WassD_p(\a,\VectMode{c})&=\left(\min_{\P\in U(\a,\VectMode{c})}\dotp{\P}{\distD^p}\right)^{1/p} \leq \dotp{\SS}{\distD^p}^{1/p}\\
&= \left(\sum_{ik} \distD^p_{ik}\sum_{j} \frac{\P_{ij}\Q_{jk}}{\tilde{\b}_j}\right)^{1/p} \leq \left(\sum_{ijk} \left(\distD_{ij}+\distD_{jk}\right)^p \frac{\P_{ij}\Q_{jk}}{\tilde{\b}_j}\right)^{1/p} \\
& \leq \left(\sum_{ijk} \distD^p_{ij} \frac{\P_{ij}\Q_{jk}}{\tilde{\b}_j}\right)^{1/p} + \left(\sum_{ijk}\distD^p_{jk} \frac{\P_{ij}\Q_{jk}}{\tilde{\b}_j}\right)^{1/p}.
\end{align*}
第1の不等式は $\SS$ の準最適性により、第2は $\distD$ の要素に対する三角不等式により、第3は Minkowski の不等式による。したがって
\begin{align*}
\WassD_p(\a,\VectMode{c})& \leq \left(\sum_{ij} \distD^p_{ij}\P_{ij} \sum_k \frac{\Q_{jk}}{\tilde{\b}_j}\right)^{1/p} + \left(\sum_{jk} \distD^p_{jk} \Q_{jk} \sum_i \frac{\P_{ij}}{\tilde{\b}_j}\right)^{1/p}\\
&= \left(\sum_{ij} \distD^p_{ij}\P_{ij}\right)^{1/p} + \left(\sum_{jk} \distD^p_{jk} \Q_{jk}\right)^{1/p}\\
&= \WassD_p(\a,\b) +\WassD_p(\b,\c),
\end{align*}
これで証明が完了する。
\end{proof}


%%%%%%%
\begin{rem}[$0<p\leq 1$ の場合]
$0 < p \leq 1$ のとき、$\distD^p$ 自体が距離であることに注意されたい。これは、$p\geq 1$ のとき $\WassD_p(\a,\b)$ が距離である一方、$p \leq 1$ の場合には実際には $\WassD_p(\a,\b)^p$ が単体上の距離を定義することを意味する。
\end{rem}
%%%%%%%


%%%%%%%
\begin{rem1}{Wasserstein 距離の応用}
OT 距離がビン間の ground 距離をそのビン上のヒストグラム間の距離に自動的に「持ち上げる」という事実は、ヒストグラムの比較のためのコンピュータビジョンと機械学習における応用にとって、OT を最適な手法とする。
%
これらの分野における古典的なアプローチは、局所特徴（例えば画像記述子）を「プーリング」し、検索、クラスタリング、分類を行うために特徴の経験分布のヒストグラム（いわゆる bag of features）を計算することである。例えば~\citep{oliva2001modeling}を参照されたい。
%
同様の考えに沿って、OT 距離は持ち上げられた特徴空間上で信号および画像解析を行うために用いることができる~\citep{thorpe2017transportation}。
%
検索とクラスタリングへの応用は画期的な論文~\citep{RubTomGui00}によって開始され、一部の応用に適した閾値行列 $\C$ に対するより高速なアルゴリズムに続く新たな応用がある。例えばコンピュータビジョンにおける~\citep{pele2008linear,Pele-iccv2009}を参照されたい。
%
より最近の応用は、bag-of-words に対する earth mover's distance の使用を強調しており、次元削減~\citep{pmlr-v51-rolet16}やテキスト分類~\citep{kusner2015word,huang2016supervised}を行うため、あるいは bag-of-words を出力する多クラス分類器を訓練するための代替損失関数を定義するため~\citep{FrognerNIPS}に用いられている。
%
\citet{kolouri2017optimal}は信号処理と機械学習へのこのような応用の最近の概観を提供している。
\end{rem1}
%%%%%%%





%%%%%%%
\begin{rem1}{測度間の Wasserstein 距離}
命題~\ref{prop-metric-histo}は、離散である必要のない任意の測度を扱うように一般化できる。

\begin{prop}\label{prop-metric-measure}
$\X=\Y$ であり、ある $p \geq 1$ に対して $\c(x,y)=\dist(x,y)^p$ であると仮定する。ここで $\dist$ は $\X$ 上の距離である。\ie
\begin{enumerate}[label=(\roman*)]
	\item  $\dist(x,y) = \dist(y,x) \geq 0$；
	\item  $\dist(x,y)=0$ となる必要十分条件は $x=y$ である；
	\item  $\foralls (x,y,z) \in \X^3, \dist(x,z) \leq \dist(x,y)+\dist(y,z)$。
\end{enumerate}
このとき $\X$ 上の $p$-Wasserstein 距離
\eql{\label{eq-defn-wass-dist}
	\Wass_p(\al,\be) \eqdef \MK_{\dist^p}(\al,\be)^{1/p}
}
（$\Wass_p$ は $\dist$ に依存することに注意）は確かに距離である。すなわち $\Wass_p$ は対称、非負であり、$\Wass_p(\al,\be)=0$ となる必要十分条件は $\al = \be$ であり、三角不等式
\eq{
	\foralls (\al,\be,\ga) \in  \Mm_+^1(\X)^3, \quad \Wass_p(\al,\ga) \leq \Wass_p(\al,\be) + \Wass_p(\be,\ga)
}
を満たす。
\end{prop}
\begin{proof}
証明は命題~\ref{prop-metric-histo}と同じアプローチに従い、$(\al,\be)$ 間と $(\be,\ga)$ 間の最適カップリングを「接着」して得られる $(\al,\ga)$ 間のカップリングの存在に依拠する。
\end{proof}
\end{rem1}
%%%%%%%


%%%%%%%
\begin{rem2}{幾何学的直観と弱収束}
Wasserstein 距離 $\Wass_p$ は多くの重要な性質を持つ。最も重要なのは、弱距離であることである。すなわち、台が重ならない特異分布（例えば離散分布）を比較し、2つの分布の台の空間的なずれを定量化することができる。
%
特に、「古典的な」距離（あるいはダイバージェンス）は離散分布の間では定義されないことさえある（$L^2$ ノルムは基底測度に関する密度を持つ連続測度にしか適用できず、離散 $\ell^2$ ノルムは位置 $(x_i,y_j)$ が適切に機能するために予め決められた離散集合の値を取る必要がある）。対照的に、任意の $p> 0$ に対して $\Wass_p^p(\de_x,\de_y) = \dist(x,y)$ が成り立つ。実際、$\Couplings(\de_x,\de_y)=\{ \delta_{x,y}\}$ であることに気づけば十分であり、したがって Kantorovich 問題は唯一の実行可能解しか持たないため、$\Wass_p^p(\de_x,\de_y)$ は必然的に $(\dist(x,y)^p)^{1/p}=\dist(x,y)$ である。これは $x \rightarrow y$ ならば $\Wass_p(\de_x,\de_y) \rightarrow 0$ であることを示す。
%
この性質は、$\Wass_p$ が以下で定義する弱収束を定量化する方法であることに対応する。

\begin{defn}[弱収束]\label{dfn-weak-conv}
	コンパクト領域 $\Xx$ 上で、$(\al_k)_k$ が $\Mm_+^1(\Xx)$ において $\al$ に弱収束する（$\al_k \rightharpoonup \al$ と表記する）必要十分条件は、任意の連続関数 $g \in \Cc(\Xx)$ に対して $\int_\Xx g \d\al_k \rightarrow \int_\Xx g \d\al$ が成り立つことである。
	%
	非コンパクト領域上では $g$ に追加的な減衰条件を課す必要がある。
	%
	この弱収束の概念はランダムベクトルの法則の収束に対応する。
\end{defn}

この収束は $\Wass_p(\al_k,\al) \rightarrow 0$ と同値であることが示せる~\citep[Theorem 6.8]{Villani09}（非有界距離空間に対しては $p$ 次までのモーメントの収束も必要）。
\end{rem2}
%%%%%%%


%%%%%%%%%%%%%%
\begin{rem1}{平行移動}
ground コスト $c(x,y)=\norm{x-y}^2$ を持つ Euclid 空間 $\X=\RR^\dim$ 上の Wasserstein 距離の良い特徴は、平行移動を因数分解できることである。実際、$T_{\tau} : x \mapsto x-\tau$ を平行移動作用素として、
\eq{
	\Wass_2(T_{\tau\sharp} \al,T_{\tau'\sharp} \be)^2  =
	\Wass_2(\al,\be)^2 - 2\dotp{\tau-\tau'}{ \mean_{\al}-\mean_{\be} } + \norm{\tau-\tau'}^2
}
が成り立つ。ここで $\mean_{\al} \eqdef \int_\X x \d\al(x) \in \RR^\dim$ は $\al$ の平均である。
%
特に、これは距離の以下のきれいな分解を意味する：
\eq{
	\Wass_2(\al,\be)^2  = \Wass_2(\tilde\al,\tilde{\be})^2  + \norm{\mean_{\al}-\mean_{\be}}^2,
}
ここで $(\tilde\al,\tilde{\be})$ は「中心化された」零平均測度 $\tilde\al=T_{\mean_{\al}\sharp}\al$ である。
\end{rem1}
%%%%%%%%%%%%%%
\pagebreak
%%%%%%%
\begin{rem1}{$p = +\infty$ の場合}\label{rem-p-inf}
非形式的に、$p \rightarrow +\infty$ における $\Wass_p^p$ の極限は
\eql{\label{eq-wass-infty}
	\Wass_{\infty}(\al,\be) \eqdef
	\umin{\pi \in \Couplings(\al,\be)}
		\usup{(x,y) \in \Supp(\pi)} d(x,y)
}
であり、ここで上限は $\X^2$ 上の測度 $\pi$ に関する本質的上限として理解されるべきである。
%
$p<+\infty$ の場合と対照的に、これは非凸最適化問題であり、数値的に解くことも理論的に研究することも困難である。$\Wass_{\infty}$ 距離は $(\al,\be)$ の台の間の Hausdorff 距離と関連している。\S~\ref{sec-hausdorff}を参照されたい。詳細については~\citep{champion2008wasserstein}を参照されたい。
\end{rem1}
%%%%%%%



%%%%%%%%%%%%%%%%%%%%%%%%%%%%%%%%%%%%%%%%%%%%
\section{双対問題}

Kantorovich 問題~\eqref{eq-mk-discr}は制約付き凸最小化問題であり、したがっていわゆる双対問題、すなわち制約付き凹最大化問題と自然に対を成す。以下の基本的な命題は主問題と双対問題の関係を説明する。

\begin{prop}\label{prop-duality-discr}
Kantorovich 問題~\eqref{eq-mk-discr}は双対
\eql{\label{eq-dual}
	\MKD_\C(\a,\b) =
	\umax{(\fD,\gD) \in \PotentialsD(\C)} \dotp{\fD}{\a} + \dotp{\gD}{\b}
}
を許容する。ここで許容双対変数の集合は
\eql{\label{eq-feasible-potential}
	\PotentialsD(\C) \eqdef \enscond{
		(\fD,\gD) \in \RR^n \times \RR^m
	}{ \foralls (i,j) \in \range{n} \times \range{m}, \fD \oplus \gD \leq \C }
}
である。このような双対変数はしばしば「Kantorovich ポテンシャル」と呼ばれる。
\end{prop}

\begin{proof}
この結果は、線形計画問題の強双対性に関するより一般的な結果~\citep[p. 148, Theo. 4.4]{bertsimas1997introduction}の直接的な帰結である。証明のより容易な部分、すなわち式~\eqref{eq-dual}の右辺が $\MKD_\C(\a,\b)$ の下界であることの確立は、次の節の注意~\ref{rem-duality}で議論される。
%
完全を期して、Lagrange 双対性を用いてこの結果を導出する。\eqref{eq-mk-discr}に付随する Lagrange 関数は
\eql{\label{eq-mk-lagr}
	\umin{\P \geq 0} \umax{ (\fD,\gD) \in \RR^n \times \RR^m }
		\dotp{\C}{\P} + \dotp{\a - \P\ones_m}{\fD} + \dotp{\b - \transp{\P} \ones_n}{\gD}
}
である。上の min と max を交換する。これは線形計画問題（有限次元）を考える場合に常に可能であり、
\eq{
	\umax{ (\fD,\gD) \in \RR^n \times \RR^m }
	\dotp{\a}{\fD} + \dotp{\b}{\gD}
	+ \umin{\P \geq 0}
		\dotp{\C - \fD\transp{\ones_m} - \ones_n \transp{\gD}}{\P}
}
が得られる。次に
\eq{
	\umin{\P \geq 0} \dotp{\Q}{\P} =
	\choice{
		0 \qifq \Q \geq 0,\\
		-\infty \quad \text{otherwise}
	}
}
であることを注意すれば、制約は $\C - \fD\transp{\ones_m} - \ones_n \transp{\gD} = \C-\fD\oplus \gD \geq 0$ と読める。
\end{proof}

Lagrange 関数~\eqref{eq-mk-lagr}の主-双対最適性関係により、最適輸送計画の台を特定できる（\S\ref{s-complementary}も参照）：
\eql{\label{eq-mk-pd-rel}
\{(i,j) \in \range{n} \times \range{m} : \P_{i,j}>0\} \subset \enscond{(i,j) \in \range{n} \times \range{m}}{ \fD_i+\gD_j=\C_{i,j} }.
}


\begin{rem}\label{rem-kantorovich-dual}
注意~\ref{rem-kantorovich-primal}で Kantorovich 問題に与えた解釈に続き、双対問題の直観的な説明を行う。その設定では、事業者は倉庫から工場へ資源をできるだけ低いコストで全体として移動させたいと考えていた。事業者は~\eqref{eq-mk-discr}を解き、$\P^\star$ に定められた指示に従い、輸送会社に $\dotp{\P^\star}{\C}$ を支払うことができる。

\emph{物流のアウトソーシング}　事業者が線形計画問題~\eqref{eq-mk-discr}を解く計算手段を持っていないとしよう。事業者はその作業をベンダーにアウトソーシングすることを決定する。ベンダーは以下の構造を持つ価格体系を選択する：ベンダーは物流作業を集荷と配送に分割し、各倉庫 $i$ で資源の1単位を集荷するのに集荷価格 $\fD_i$ を適用し（その単位がどこに送られるかに関わらず）、工場 $j$ に資源の1単位を配送するのに価格 $\gD_j$ を適用する（その単位がどの倉庫から来るかに関わらず）。倉庫 $i$ にはちょうど $\a_i$ 単位があり、工場 $j$ には $\b_j$ が必要であるため、ベンダーはこの価格体系の帰結として事業者の物流問題を解決するのに $\dotp{\fD}{\a}+\dotp{\gD}{\b}$ の価格を要求する。

\emph{価格の設定}　ベンダーが用いる価格体系は、かなり自然に任意の負の価格を許容することに注意されたい。実際、ベンダーが倉庫に対して価格ベクトル $\fD$、工場に対して価格ベクトル $\gD$ を適用する場合、$\fD$ のすべての成分を任意の数だけ同時に減少させ、$\gD$ のすべての成分を同じ数だけ増加させても、合計請求額は変わらない。すべての倉庫の資源の総量はすべての工場に配送されなければならない量に等しいためである。言い換えれば、ベンダーは事業者の商品の一部を集荷するのに支払うことで極めて良い取引の錯覚を与え、配送にさらに多く請求することで単にその損失を補うことができる。これを知ったベンダーは、そのサービスにできるだけ多く請求したいと考え、ベクトル $\fD$ と $\gD$ をできるだけ高く設定する。

\emph{価格の確認}　競合するベンダーがいない場合、事業者はベンダーの価格が合理的であるかを確認する迅速な方法を考えなければならない。一つの可能な方法は、事業者が問題~\eqref{eq-mk-discr}を解いて最も効率的な計画の価格 $\MKD_{\C}(\a,\b)$ を計算し、ベンダーの提示がその金額以下であるかを確認することである。しかし、事業者はそもそもそのような長い計算を行う余裕がないことを思い出されたい。幸い、ベンダーが競争力のある提示をしているかを事業者が確認するためのはるかに効率的な方法がある。$\fD_i$ はベンダーが $i$ で1単位を集荷するのに請求する価格であり、$\gD_j$ は $j$ で1単位を配送する価格であることを思い出されたい。したがって、ベンダーの価格体系は $i$ から $j$ へ資源の1単位を移送するコストが正確に $\fD_i+\gD_j$ であることを意味する。しかし、事業者は輸送会社が $i$ から $j$ へ1単位を輸送するのに $\C_{i,j}$ と価格設定していることも知っている。したがって、任意のペア $i,j$ に対して合計価格 $\fD_i+\gD_j$ が $\C_{i,j}$ より厳密に大きい場合、ベンダーはその作業に対して輸送会社が請求する公正な価格を超える請求をしており、事業者はベンダーの提示を拒否すべきである。


\begin{figure}[h!]
\centering
\includegraphics[width=\linewidth]{dual-solutions-overlay/dual_sols_final}
\caption{\label{dual-solutions-overlay}
左のプロットでは、それぞれ青い点と赤い四角で表される2つの離散測度 $\al$ と $\be$ の間の optimal transport 問題を考える。これらのマーカーの面積は各位置の重みに比例する。このプロットは二乗 Euclid コストを用いた最適輸送 $\P^\star$ も表示している。対応する双対（Kantorovich）ポテンシャル $\fD^\star$ と $\gD^\star$ は右のプロットにも表示されている。$\alpha$ の各点に「価格」$\fD_i^\star$ があるため（$\gD$ と $\beta$ についても同様）、その点の色は右のカラーマップを用いて得られた値を表す。これらのポテンシャルは相対価格として解釈できる。すなわち、最善の輸送計画の下で、$\al$ の各位置から質量を移動させる個別コスト、あるいは逆に $\be$ の任意の点に質量を送り込むコストを示す。したがって、最適輸送コストは、左の図のすべての弧の二乗長を太さで重み付けした和、あるいは代替的に双対定式化を用いて、右のプロットの各マーカーの面積に値（色で符号化）を乗じた和に等しい。}
\end{figure}


\emph{双対問題としての最適価格}　したがって、事業者はすべてのペア $i,j$ に対してベンダーが提示する価格が $\fD_i+\gD_j \leq \C_{i,j}$ を満たすことを確認することが利益になる。事業者はベンダーがこれら $n\times m$ 個の不等式に適合する価格ベクトルを提供したことを確認するとしよう。事業者はベンダーの提案が魅力的であると結論づけられるだろうか？概算を行い、事業者は確かにその提示を受け入れることが利益であると結論づける。実際、事業者の任意の輸送計画 $\P$ のコストは $\dotp{\P}{\C}=\sum_{i,j} \P_{i,j}\C_{i,j}$ であるため、事業者はこれら $n\times m$ 個の不等式を適用して、任意の輸送計画 $\P$（最適なもの $\P^\star$ を含む）に対して、周辺制約により
\eq{\begin{aligned}
	\sum_{i,j} \P_{i,j}\C_{i,j} &\geq \sum_{i,j} \P_{i,j} \left(\fD_{i} + \gD_j\right) = \left(\sum_{i} \fD_{i} \sum_{j} \P_{i,j}\right) + \left(\sum_{j} \gD_{j} \sum_{i} \P_{i,j}\right)\\ 	& =\dotp{\fD}{\a}+\dotp{\gD}{\b}
\end{aligned}}
が成り立つことを結論づけ、自分で作業を行う\emph{あらゆる}試みは必然的にベンダーの価格よりも高くつくことを観察する。

これを知ったベンダーは、$\dotp{\fD}{\a}+\dotp{\gD}{\b}$ を最大化する価格 $\fD,\gD$ の集合を見つけなければならないが、提示が受け入れられるためにはすべての $i,j$ に対して少なくとも基本的な不等式 $\fD_i+\gD_j\leq \C_{i,j}$ を満たさなければならない。これが問題~\eqref{eq-dual}に帰結する。後に\S\ref{s-kantorovich}で示すように、ベンダーが得る最良の価格は、事業者が $\MKD_{\C}(\a,\b)$ を計算して得る最良のコストと実際に正確に等しい。

図~\ref{dual-solutions-overlay}は同じ輸送問題から得られる主問題と双対問題の解を示す。左側では、青い点が倉庫を、赤い点が工場を表す。これらの点の面積は確率重み $\a,\b$ を表し、それらの間のリンクは最適輸送を表し、その幅は移送量に比例する。ベンダーが問題~\eqref{eq-dual}を最適化した結果として得られる最適価格が右に示されている。価格はその平均が 0 に等しくなるように選ばれている。最も高い相対価格は、図の左下にある孤立した倉庫での物資の集荷と、右上の領域にある工場への物資の配送に起因する。
\end{rem}



%%%%%%%
\begin{rem2}{任意の測度間の双対問題}
	この主-双対構成を任意の測度に拡張するためには、測度は連続関数と双対性の関係にあること（測度は連続関数に対する積分を通じてのみアクセスできる）を認識することが重要である。この双対性は以下の命題で形式化され、離散測度を扱う場合は命題~\ref{prop-duality-discr}に帰着する。

	\begin{prop}
	次が成り立つ：
	\eql{\label{eq-dual-generic}
		\MK_\c(\al,\be) =
		\usup{(\f,\g) \in \Potentials(\c)}
			\int_\X \f(x) \d\al(x) + \int_\Y \g(y) \d\be(y),
	}
	ここで許容双対ポテンシャルの集合は
	\eql{\label{eq-dfn-pot-dual}
		\Potentials(\c) \eqdef \enscond{
			(\f,\g) \in \Cc(\X) \times \Cc(\Y)
		}{
			\forall (x,y), \f(x)+\g(y) \leq \c(x,y)
		}
	}
	である。ここで $(\f,\g)$ は連続関数の対であり、離散的な場合と同様に「Kantorovich ポテンシャル」とも呼ばれる。
	\end{prop}

	離散の場合~\eqref{eq-dual}は双対ベクトルが連続ポテンシャルのサンプルであることに対応する。\ie $(\fD_i,\gD_j)=(\f(x_i),\g(y_j))$。
	%
	主-双対最適性条件により最適計画の台を追跡でき、\eqref{eq-mk-pd-rel}は
	\eql{\label{eq-mk-pd-rel-cont}
		\Supp(\pi) \subset \enscond{(x,y) \in \X \times \Y}{ \f(x)+\g(y)=\c(x,y) }
	}
	と一般化される。

	主問題~\eqref{eq-mk-generic}とは対照的に、\eqref{eq-dual-generic}の解の存在を示すことは自明ではないことに注意されたい。制約集合 $\Potentials(\c)$ がコンパクトでなく、最小化すべき関数が強圧的でないためである。
	%
	\S\ref{s-c-transform}で詳述する $c$-変換の手法を用いれば、$c(x,y)=\dist(x,y)^p$（$p\geq 1$）の場合に、最適な $(\f,\g)$ が必然的に Lipschitz 正則であることを示すことができ、これにより制約をコンパクトなものに置き換えることが可能になる。
\end{rem2}
%%%%%%%

%%%%%%%
\begin{rem2}{制約なし双対}\label{rem-uncons-dual}
$\int_\X \d\mu=\int_\Y\d\nu = 1$ の場合、制約付き双対問題~\eqref{eq-dual-generic}は制約なしの問題に置き換えることができる：
\eql{\label{eq-dual-uncons}
	\MK_\c(\al,\be) =
		\usup{(\f,\g) \in \Cc(\X) \times \Cc(\Y)}
			\int_\X \f\d\al + \int_\Y \g \d\be
			+ \umin{\X \otimes \Y} ( c - f \oplus g ),
}
ここで $(f\oplus g)(x,y)=f(x)+g(y)$ と表記した。ここでの最小値は測度 $\mu \otimes \nu$ に付随する本質的上限と見なすべきであり、\ie $f$ や $g$ が $\mu$ および $\nu$ の零測度集合上で変更されても変わらない。
%
この代替的な双対定式化は Francis Bach によって我々に指摘されたものである。
%
これは主問題~\eqref{eq-mk-generic}に冗長な制約 $\int \d\pi=1$ を加えることで得られる。
\end{rem2}
%%%%%%%


%%%%%%%
\begin{rem2}{Monge--Kantorovich の同値性 --- Brenier の定理}\label{rem-exist-mongemap}
以下の定理はしばしば~\citet{Brenier91}に帰せられ、$\RR^\dim$ において $p=2$ の場合に、2つの入力測度の少なくとも一方が密度を持ち、2次モーメントを持つ測度に対して、Kantorovich 問題と Monge 問題が同値であることを保証する。興味のある読者は、ほぼ同時期に出版された同じ結果の変形版~\citep{cuesta1989,ruschendorf1990characterization}、特に~\citep{MR923203}における元の結果、および~\citet{knott1984optimal}による先駆的研究も参照されたい。

\begin{thm}[Brenier の定理]
	$\X=\Y=\RR^\dim$ かつ $c(x,y)=\norm{x-y}^2$ の場合に、2つの入力測度の少なくとも一方（$\al$ とする）が Lebesgue 測度に関する密度 $\density{\al}$ を持つならば、Kantorovich 定式化~\eqref{eq-mk-generic}における最適な $\pi$ は一意であり、「Monge 写像」$\T : \RR^\dim \rightarrow \RR^\dim$ のグラフ $(x,\T(x))$ 上に台を持つ。すなわち $\pi = (\Id,\T)_{\sharp} \mu$ であり、\ie
\eql{\label{eq-brenier-map}
	\foralls h \in \Cc(\X \times \Y), \quad
		\int_{\X \times \Y} h(x,y) \d \pi(x,y) = \int_{\X} h(x,\T(x)) \d\mu(x).
}
さらに、この写像 $\T$ は凸関数 $\phi$ の勾配 $\T(x) = \nabla\phi(x)$ として一意的に定義される。ここで $\phi$ は $(\nabla\phi)_\sharp \mu=\nu$ を満たす（加法定数を除いて）一意的な凸関数である。この凸関数は~\eqref{eq-dual-generic}を解く双対ポテンシャル $\f$ と $\phi(x)=\frac{\norm{x}^2}{2}-\f(x)$ の関係にある。
\end{thm}

\begin{proof}
	証明の主要な要素の概略を述べる。詳細は例えば~\citep{SantambrogioBook}にある。
	%
	$\int c \d\pi = C_{\al,\be} - 2 \int \dotp{x}{y}\d\pi(x,y)$ であることに注意する。ここで定数は $C_{\al,\be} = \int \norm{x}^2\d\al(x) + \int \norm{y}^2\d\be(y)$ である。\eqref{eq-mk-generic}を解く代わりに、問題
	\eq{
		\umax{\pi \in \Couplings(\al,\be)}
			\int_{\X \times \Y} \dotp{x}{y} \d\pi(x,y)
	}
	を考えることができる。その双対は
	\eql{\label{eq-brenier-proof-1}
		\umin{(\phi,\psi)}
		\enscond{
			\int_\X \phi \d\al + \int_\Y \psi \d\be
		}{
			\forall (x,y), \quad \phi(x)+\psi(y) \geq \dotp{x}{y}
		}
	}
	となる。これらの変数と~\eqref{eq-dfn-pot-dual}のそれとの関係は
	$(\phi,\psi) = (\frac{\norm{\cdot}^2}{2}-\f,\frac{\norm{\cdot}^2}{2}-\g)$ である。
	%
	制約は
	\eql{\label{eq-brenier-proof-leg}
		\foralls y, \quad \psi(y) \geq \phi^*(y) \eqdef \usup{x} \dotp{x}{y} - \phi(x)
	}
	で置き換えることができる。ここで $\phi^*$ は $\phi$ の Legendre 変換であり、線形形式の上限として凸関数である（\eqref{eq-legendre}も参照）。\eqref{eq-brenier-proof-2}に現れる目的関数が線形であり積分する測度が正であるため、$\psi$ に関して明示的に最小化でき、$\psi=\phi^*$ として制約なし問題
	\eql{\label{eq-brenier-proof-2}
		\umin{\phi} \int_\X \phi \d\al + \int_\Y \phi^* \d\be
	}
	を考えることができる。\S\ref{sec-c-transforms}および\S\ref{s-c-transform}も参照されたい。そこではこの考え方がそれぞれ離散的な設定と一般コスト $\c(x,y)$ に対して適用される。
	%
	この議論を2回繰り返すことで、$\phi$ を $\phi^{**}$（凸関数）で置き換えることができ、したがって~\eqref{eq-brenier-proof-2}において $\phi$ が凸であることを課すことができる。
	%
	条件~\eqref{eq-mk-pd-rel-cont}は、最適な $\pi$ が $\enscond{(x,y)}{\phi(x)+\phi^*(y)=\dotp{x}{y}}$ 上に台を持つことを示し、このような $y$ は Legendre 変換の最小化~\eqref{eq-brenier-proof-leg}で最適であり、その最適性条件は $y \in \partial \phi(x)$ と読める。
	%
	$\phi$ は凸であるため、ほとんど至る所で微分可能であり、$\al$ が密度を持つため、$\al$-ほとんど至る所でも微分可能である。
	%
	これは各 $x$ に対して、対応する $y$ が $\al$-ほとんど至る所で $y = \nabla\phi(x)$ として一意に定義されることを示し、必然的に $\pi = (\Id,\nabla\phi)_\sharp \al$ であることを示す。
\end{proof}

この結果は、非特異密度を持つ $\Wass_2$ の設定において、Monge 問題~\eqref{eq-monge-continuous}とその Kantorovich 緩和~\eqref{eq-mk-generic}が等しいこと（緩和がタイトであること）を示している。これは、割当ての場合の命題~\ref{prop-matching-kanto}の連続版であり、周辺が等しく一様のとき optimal transport 問題の最小値が置換行列（離散写像）で達成されることを述べている。
%
Brenier の定理は、最適輸送写像が凸関数の勾配でなければならないと述べることで、1次元より大きい次元における増加関数の概念の有用な一般化を提供する。これが、optimal transport が任意の次元における分位関数の定義に用いられる主な理由であり、分位回帰問題への応用に有用である~\citep{carlier2016vector}。

この定理は多くの方向に拡張できることにも注意されたい。
%
$\al$ が密度を持つという条件は、Hausdorff 次元が $\dim-1$ より小さい「小さな集合」（\emph{例えば}超曲面）に質量を与えないという条件に弱めることができる。
%
$\c(x,y)=h(x-y)$ の形のコストも考慮でき、ここで $h$ は狭義凸関数である。
\end{rem2}
%%%%%%%


%%%%%%%
\begin{rem2}{Monge--Amp\`ere 方程式}\label{rem:MA}
密度を持つ測度に対して、\eqref{eq-pfwd-density}を用いると、$\phi$ は以下の Monge--Amp\`ere 型方程式を解く（加法定数を除いて）一意的な凸関数であることが得られる：
\eql{\label{eq-monge-ampere}
	\det(\partial^2\phi(x))  \density{\be}(\nabla\phi(x)) = \density{\al}(x),
}
ここで $\partial^2\phi(x) \in \RR^{\dim \times \dim}$ は $\phi$ の Hesse 行列である。Monge--Amp\`ere 作用素 $\det(\partial^2\phi(x))$ は非線形退化 Laplacian として理解できる。小さな変位の極限 $\phi=\Id + \epsilon\psi$ では、滑らかな写像に対して
\eq{
	\det(\partial^2\phi(x)) = 1 + \epsilon \Delta \psi(x) + o(\epsilon)
}
であるため、実際に線形化として Laplacian $\Delta$ が回復される。
%
凸性制約は $\det(\partial^2\phi(x)) \geq 0$ を強制し、この方程式が解を持つために必要である。
%
Monge--Amp\`ere 方程式の理論的解析、特にその解の正則性に関する膨大な文献がある。例えば~\citep{gutierrez2016monge}を参照されたい。レビュー論文としては~\citet{caffarelli2003monge}を参照されたい。
%
主要な困難は、一般には解が滑らかである必要がなく、入力測度が任意（\emph{例えば} Dirac 測度）の場合には Alexandrov 解の手法に頼らなければならないことである。
%
固定された右辺 $f$ を持つ Monge--Amp\`ere 方程式 $\det(\partial^2 \phi(x)) = f(x)$ のより単純な場合には、多くの解法が提案されている。例えば~\citep{benamou2016monotone}およびそこでの参考文献を参照されたい。特に、異方性凸関数の捕捉には特別な注意が必要であり、通常の有限差分は不正確になりうる。
%
optimal transport に対しては、$f$ が実際には $\nabla \phi$ に依存するため、方程式~\eqref{eq-monge-ampere}の離散化と境界条件は~\citep{benamou2014numerical}およびそこでの参考文献で概説される技術的課題を生じさせる。
%
また、不動点反復に基づく関連する解法が画像レジストレーションに適用されていることにも注意されたい~\citep{haker2004optimal}。
\end{rem2}
%%%%%%%


%%%%%%%%%%%%%%%%%%%%%%%%%%%%%%%%%%%%%%%%%%%%
\section{特殊な場合}\label{sec:specialcases}

一般に、OT 距離の計算は数値的に手間がかかる。\S\S\ref{c-algo-basics}、\ref{c-entropic}、\ref{c-dynamic}で異なる数値解法を詳述する前に、まず OT 問題の解決が比較的容易な特殊な好ましい場合をレビューする。


%%%%%%%
\begin{rem}[二値コスト行列と1ノルム]\label{rem-binary}
コスト行列 $\C$ が対角上で 0 でそれ以外で $1$ のとき、すなわち $\C=\ones_{n\times n}-\Identity_n$ のとき、$\a$ と $\b$ の間の 1-Wasserstein 距離はそれらの差の1ノルムに等しいことが容易に確認できる：$\MKD_\C(\a,\b)=\norm{\a-\b}_1$。
\end{rem}
%%%%%%%

%%%%%%%
\begin{rem1}{Kronecker コスト関数と全変動}
上の注意~\ref{rem-binary}に加えて、$c(x,y)$ が $x=y$ のとき $0$ で $x\ne y$ のとき $1$ である場合にこの結果が任意の測度に拡張されることも容易に確認できる。2つの離散測度 $\al$ と $\be$ の間の OT 距離はそれらの全変動距離に等しい（例~\ref{exmp-tv}も参照）。
\end{rem1}
%%%%%%%



%%%%%%%
\begin{rem1}{1次元の場合 --- 経験測度}\label{rem-1d-empir}
ここで $\X=\RR$ とする。$\al = \frac{1}{n}\sum_{i=1}^n \de_{x_i}$ かつ $\be = \frac{1}{n}\sum_{j=1}^n \de_{y_j}$ と仮定し、（一般性を失うことなく）点が順序付けられている、\ie $x_1 \leq x_2 \leq \cdots \leq x_n$ かつ $y_1 \leq y_2 \leq \cdots \leq y_n$ と仮定すると、単純な公式
\eql{\label{eq-1d-empirical}
	\Wass_p(\al,\be)^p = \frac{1}{n} \sum_{i=1}^n |x_i-y_i|^p
}
が成り立つ。\ie 局所的には（点がすべて異なると仮定すれば）、$\Wass_p(\al,\be)$ は $\al$ と $\be$ の順序付けられた値の2つのベクトル間の $\ell^p$ ノルムである。この主張は局所的にのみ有効であり、値のいくつかが変化するたびに順序（およびそれらのベクトル表現）が変わりうるという意味においてである。この公式は注意~\ref{rem-1d-ot-generic}で詳述されるより一般的な設定の単純な帰結である。
%
図~\ref{fig-1d-discrete}の上段は、同じ点数を持つ経験測度間の1次元輸送写像を示す。下段は、この単調写像が任意の離散測度にどのように一般化されるかを示す。


円上の OT を効率的に計算するためにこの1次元計算を利用することも可能である。\citet{delon-circle}によって示されているように。
%
コストが距離の凹関数である場合、特に $p<1$ のとき、最適輸送計画の挙動は大きく異なるが、効率的な解法も存在する~\citep{delon-concave}。
\end{rem1}
%%%%%%%


\begin{figure}[h!]
\centering
\includegraphics[width=.65\linewidth]{1d-discrete/1d-schematic}
\caption{\label{fig-1d-discrete}
1次元最適カップリング：各矢印 $x_i \rightarrow y_j$ は最適カップリングにおける非零の $\P_{i,j}$ を示す。
%
上段：同じ点数を持つ経験測度（最適マッチング）。下段：一般的な場合。
%
これは単調再配置に対応する。$x_i \leq x_{i'}$ が $\P_{i,j} \neq 0, \P_{i',j'} \neq 0$ を満たすならば、必然的に $y_j \leq y_{j'}$ である。
}
\end{figure}


%%%%%%%
\begin{rem1}{ヒストグラム均等化}
1次元 optimal transport はヒストグラム均等化の実行に用いることができ、グレースケール画像のパレットの正規化に応用される。図~\ref{fig-hist-eq}を参照されたい。この場合、$(\bar x_i)_i$ と $(\bar y_j)_j$ は2つの入力画像のすべてのピクセルのグレー色レベル（黒は $0$、白は $1$、その間のすべての値）を所定の順序（\ie 列方向）で列挙したものを表す。各画像のピクセル数が同じで $n\times m$ に等しいと仮定すると、これらの色レベルをソートすることで注意~\ref{rem-1d-empir}のように $x_i = \bar x_{\si_1(i)}$ と $y_j = \bar y_{\si_2(j)}$ が定義される。ここで $\si_1, \si_2 : \{1,\ldots,nm\} \rightarrow \{1,\ldots,nm\}$ は置換であり、$\si \eqdef \si_2 \circ \si_1^{-1}$ は2つの離散分布間の最適割当てである。画像処理の応用では、$(\bar y_{\si(i)})_i$ は $\bar x$ の均等化されたバージョンの色値を定義し、その経験分布は $\bar y$ のそれと正確に一致する。その画像の均等化バージョンは、この $nm$ 次元ベクトルをサイズ $n\times m$ の画像として折り返すことで復元できる。また、$t \in [0,1] \mapsto (1-t) \bar x_{i} + t \bar y_{\si(i)}$ は元画像と均等化画像の間の補間を定義し、そのピクセルの経験分布は入力のそれらの間の変位補間（\eqref{eq-displacement-1d-cumul}で定義される通り）である。
\end{rem1}
%%%%%%%

\begin{figure}[h!]
\centering
\begin{tabular}{@{}c@{\hspace{1mm}}c@{\hspace{1mm}}c@{\hspace{1mm}}c@{\hspace{1mm}}c@{\hspace{1mm}}@{}}
\includegraphics[width=.19\linewidth]{hist-eq/equalized-1}&
\includegraphics[width=.19\linewidth]{hist-eq/equalized-3}&
\includegraphics[width=.19\linewidth]{hist-eq/equalized-5}&
\includegraphics[width=.19\linewidth]{hist-eq/equalized-7}&
\includegraphics[width=.19\linewidth]{hist-eq/equalized-9}\\
\includegraphics[width=.19\linewidth]{hist-eq/equalize-1}&
\includegraphics[width=.19\linewidth]{hist-eq/equalize-3}&
\includegraphics[width=.19\linewidth]{hist-eq/equalize-5}&
\includegraphics[width=.19\linewidth]{hist-eq/equalize-7}&
\includegraphics[width=.19\linewidth]{hist-eq/equalize-9}\\
$t=0$ & $t=0.25$& $t=0.5$& $t=.75$& $t=1$
\end{tabular}
\caption{\label{fig-hist-eq}
画像処理のためのヒストグラム均等化。$t$ はヒストグラム間の変位補間をパラメータ化する。
}
\end{figure}


%%%%%%%
\begin{rem2}{1次元の場合 --- 一般の場合}\label{rem-1d-ot-generic}
$\RR$ 上の測度 $\al$ に対して、$\RR$ から $[0,1]$ への累積分布関数を
\eql{\label{eq-cumul-defn}
	\foralls x \in \RR, \quad \cumul{\al}(x) \eqdef \int_{-\infty}^x \d\al
}
と導入し、その擬似逆関数 $\cumul{\al}^{-1} : [0,1] \rightarrow \RR \cup \{-\infty\}$ を
\eql{\label{eq-pseudo-inv-cum}
	\foralls r \in [0,1], \quad \cumul{\al}^{-1}(r) = \umin{x} \enscond{x \in \RR \cup \{-\infty\} }{ \cumul{\al}(x) \geq r }
}
と導入する。
%
この関数は $\alpha$ の一般化分位関数とも呼ばれる。任意の $p \geq 1$ に対して
\eql{\label{eq-wass-cumul}
	\Wass_p(\al,\be)^p = \norm{ \cumul{\al}^{-1} - \cumul{\be}^{-1} }_{L^p([0,1])}^p = \int_0^1 | \cumul{\al}^{-1}(r) - \cumul{\be}^{-1}(r) |^p \d r
}
が成り立つ。これは、写像 $\al \mapsto \cumul{\al}^{-1}$ を通じて、Wasserstein 距離が $L^p$ ノルムを備えた線形空間に等距離的であること、あるいは同値に、実数直線上の測度に対する Wasserstein 距離が Hilbert 的距離であることを意味する。
これは1次元 optimal transport の幾何学を非常に単純にするが、命題~\ref{prop-negative-definite}で議論され、より一般的には\S\ref{sec-non-embeddability}で論じられるように、Hilbert 的でない高次元の幾何学とは大きく異なる。
%
$p=1$ に対しては、さらに単純な公式
\begin{align}\label{eq-w1-1d}
	\Wass_1(\al,\be) &= \norm{ \cumul{\al} - \cumul{\be} }_{L^1(\RR)} =
	\int_\RR | \cumul{\al}(x) - \cumul{\be}(x) | \d x \\
	&= \int_\RR \abs{ \int_{-\infty}^x \d(\al-\be) } \d x
\end{align}
が成り立ち、$\Wass_1$ がノルムであることを示す（任意の次元への一般化は\S\ref{sec-w1-eucl}を参照）。
%
最適 Monge 写像 $\T$（$\T_\sharp \al=\be$ を満たす）は
\eql{\label{eq-OT-map-1d}
 	\T = \cumul{\be}^{-1} \circ \cumul{\al}
}
で定義される。図~\ref{fig-1d-ot}は累積関数を通じた1次元 OT の計算を示す。また、\eqref{eq-displacement-1d-cumul}で詳述されるように計算された変位補間も表示する。注意~\ref{rem-bary-1d}も参照されたい。1次元における optimal transport の性質の詳細な調査については、\cite[Chapter 2]{SantambrogioBook}を参照されたい。
\end{rem2}
%%%%%%%


\newcommand{\MyFigCumulMeas}[1]{\includegraphics[width=.33\linewidth]{1d-cumulative/#1}}
\newcommand{\MyFigCumulCum}[1]{\includegraphics[width=.24\linewidth]{1d-cumulative/#1}}
\begin{figure}[ht!]
\centering
\begin{tabular}{@{}c@{\hspace{1mm}}c@{\hspace{1mm}}c@{}}
\MyFigCumulMeas{input-mu}&
\MyFigCumulMeas{input-nu}&
\MyFigCumulMeas{interp-bary}\\
$\mu$ & $\nu$ & ${ (t\T+(1-t)\Id)_\sharp \mu}$
\end{tabular}
%%%%
\begin{tabular}{@{}c@{\hspace{2mm}}c@{\hspace{2mm}}c@{\hspace{2mm}}c@{}}
\MyFigCumulCum{cumul}&
\MyFigCumulCum{icumul}&
\MyFigCumulCum{transports}&
\MyFigCumulCum{interp-cumul}\\
$(\cumul{\al},\cumul{\be})$ &
$(\cumul{\al}^{-1},\cumul{\be}^{-1})$ &
$(T,T^{-1})$ &
$(1-t)\cumul{\al}^{-1}+t\cumul{\be}^{-1}$
\end{tabular}
\caption{\label{fig-1d-ot}
\eqref{eq-OT-map-1d}で詳述された累積関数を用いた、2つの1次元測度間の OT と変位補間の計算。
}
\end{figure}

%%%%%%%
\begin{rem2}{Gauss 分布間の距離}\label{rem-dist-gaussians}
$\al = \Nn(\mean_\al,\cov_\al)$ と $\be = \Nn(\mean_\be,\cov_\be)$ が $\RR^\dim$ の2つの Gauss 分布であるとき、以下の写像
\eql{\label{eq-transport-Bures}T:x\mapsto \mean_\be + A(x-\mean_\al),}
ここで
$$A=\cov_\al^{-\tfrac{1}{2}}\Big(\cov_\al^{\tfrac{1}{2}}\cov_\be\cov_\al^{\tfrac{1}{2}}\Big)^{\tfrac{1}{2}}\cov_\al^{-\tfrac{1}{2}}=\transp{A},$$
が $T_\sharp \rho_\al = \rho_\be$ を満たすことが示される。実際、変数変換公式~\eqref{eq-pfwd-density}が以下のように満たされることに気づけばよい：
$$
\begin{aligned}\rho_\be(T(x))&=\det(2\pi\cov_\be)^{-\tfrac{1}{2}} \exp(-\dotp{T(x)-\mean_\be}{\cov_\be^{-1}(T(x)-\mean_\be)})\\
&= \det(2\pi\cov_\be)^{-\tfrac{1}{2}} \exp(-\dotp{ x-\mean_\al}{\transp{A}\cov_\be^{-1}A(x-\mean_\al)}) \\
&= \det(2\pi\cov_\be)^{-\tfrac{1}{2}} \exp(-\dotp{ x-\mean_\al}{\cov_\al^{-1}(x-\mean_\al)}),
\end{aligned}$$
また $T$ は線形写像であるため
$$|\det T'(x)|= \det A = \left(\frac{\det\cov_\be}{\det\cov_\al}\right)^{\tfrac{1}{2}}$$
であり、したがって $\rho_\al=|\det T'| \rho_\be$、すなわち $T_\sharp \al = \be$ が回復される。$T$ は凸関数 $\psi:x\mapsto \tfrac{1}{2}\dotp{x-\mean_\al}{A (x-\mean_\al)} + \dotp{\mean_\be}{x}$ の勾配であることに気づき、\citeauthor{Brenier91}の定理~\citeyearpar{Brenier91}（注意~\ref{rem-exist-mongemap}を参照）を用いて $T$ が最適であると結論づける。その写像 $T$ と対応するポテンシャル $\psi$ は図~\ref{fig-gaussians-2d-T}および図~\ref{fig-gaussians-2d-psi}に示されている。

$\rho_\al$ の1次および2次モーメントを含む追加的な計算により、その写像の輸送コストは
\eql{\label{eq-dist-gauss}
	\Wass_2^2( \al,\be ) = \norm{ \mean_\al - \mean_\be }^2 + \Bb(\cov_\al,\cov_\be)^2
}
であることが得られる。ここで $\Bb$ はいわゆる \citeauthor{bures1969extension} 距離~\citeyearpar{bures1969extension}であり（\citet{,forrester2016relating}も参照）、正定値行列間の距離で
\eql{\label{eq-bure-defn}
	\Bb(\cov_\al,\cov_\be)^2 \eqdef \tr\pa{
		\cov_\al + \cov_\be - 2 ( \cov_\al^{1/2} \cov_\be \cov_\al^{1/2} )^{1/2}
	}
}
である。ここで $\cov^{1/2}$ は行列平方根である。$\Bb$ が共分散行列上の距離であること、および $\Bb^2$ が両方の引数に関して凸であることが示される。$\cov_\al = \diag(r_i)_i$ と $\cov_\be = \diag(s_i)_i$ が対角行列の場合、Bures 距離は Hellinger 距離
\eq{
	\Bb(\cov_\al,\cov_\be) = \norm{ \sqrt{r}-\sqrt{s} }_2
}
となる。1次元 Gauss 分布に対して、$\Wass_2$ は Gauss 分布の平均と標準偏差 $(\mean,\sqrt{\cov})$ をプロットする2次元平面上の Euclid 距離であり、図~\ref{fig-1d-gaussian}に示されている。
%
Gauss 分布の Wasserstein 幾何学の詳細な取り扱いについては~\citet{takatsu2011wasserstein}を参照されたい。Bures 距離に関する追加的な考察については最近の文献~\citep{malago2018wasserstein,bhatia2018bures}を参照されたい。単語の確率的埋め込みを計算するためのこの距離の最近の応用については~\citep{NIPS2018_8226}を、Kalman フィルタリングのロバストな拡張を計算するための使用については~\citep{NIPS2018_8067}を、再生核 Hilbert 空間における共分散関数に適用したものについては~\citep{NIPS2017_7149}を参照されたい。
\end{rem2}
%%%%%%
\begin{figure}[h!]
\centering
\includegraphics[width=.7\linewidth]{gaussians-2d/gaussians1.pdf}
\caption{\label{fig-gaussians-2d-T} 2つの Gauss 分布 $\rho_\al$ と $\rho_\be$。それぞれの密度の等高線で表される。平均と分散行列はそれぞれ $\mean_\al=(-2,0),\cov_\al=\frac{1}{2}\left(1 -\tfrac{1}{2};-\tfrac{1}{2}  1\right)$ と $\mean_\be=(3,1), \cov_\be=\left(2, \tfrac{1}{2}; \tfrac{1}{2}, 1\right)$ である。矢印は平面上のランダムな点 $x$ から始まり、対応する写像先 $T(x)=\mean_\be + A(x-\mean_\al)$ で終わる。}
\end{figure}



%%%%%%%
\begin{rem2}{楕円等高分布間の距離}\label{rem-dist-elliptic}
\citeauthor{gelbrich1990formula}は注意~\ref{rem-dist-gaussians}で提供されたものより一般的な結果を提供する：Gauss 分布間の Bures 距離はより一般的に\emph{楕円等高分布}に拡張される~\citeyearpar{gelbrich1990formula}。
%
要約すると、まず、与えられた平均と共分散行列を持つ2つの測度に対して、これらのパラメータを持つ2つの Gauss 分布間の距離が2つの測度間の Wasserstein 距離の下界であることが示される~\citep[Theorem 2.1]{gelbrich1990formula}。加えて、閉形式~\eqref{eq-dist-gauss}は楕円等高密度の族に拡張される：2つの密度 $\rho_\al$ と $\rho_\be$ がそのような族に属する場合、すなわち $\rho_\al$ と $\rho_\be$ が任意の点 $x$ に対して平均と正定値パラメータを用いて

\eq{\begin{aligned}\rho_\al(x) = \frac{1}{\sqrt{\det(\A)}}h(\dotp{x-\mean_\al}{\A^{-1}(x-\mean_\al)})\\ \rho_\be(x)= \frac{1}{\sqrt{\det(\B)}} h(\dotp{x-\mean_\be}{\B^{-1}(x-\mean_\be)})\end{aligned}}
と書ける場合（同じ非負値関数 $h$ で積分 $$\int_{\RR^\dim} h(\dotp{x}{x})dx=1$$ を満たすもの）、それらの最適輸送写像もまた線形写像\eqref{eq-transport-Bures}であり、Wasserstein 距離も式$\eqref{eq-dist-gauss}$で与えられる。ただし、Bures 距離のスケーリングが生成関数 $h$ にのみ依存してわずかに異なる。例えば、このスケーリングは Gauss 分布（$h(t)=e^{-t/2}$）では $1$ であり、楕円体上の一様分布（$h$ は $[0,1]$ の指示関数）では $1/(\dim+2)$ である。この結果は、楕円分布の共分散行列がその正定値パラメータの定数倍であるという事実~\citep[Theo. 4(ii)]{gomez2003survey}と、楕円分布間の Wasserstein 距離がそれらの共分散行列間の Bures 距離の関数であるという事実~\citep[Cor. 2.5]{gelbrich1990formula}から従う。
\end{rem2}
%%%%%%
\begin{figure}[h!]
\centering
\includegraphics[width=.85\linewidth]{gaussians-2d/gaussians2.pdf}
\caption{\protect\label{fig-gaussians-2d-psi} 図~\ref{fig-gaussians-2d-T}で定義されたものと同じ Gauss 分布 $\rho_\al$ と $\rho_\be$。今回は曲面として表現されている。上方の曲面は $T=\nabla \psi$ を満たす Brenier ポテンシャル $\psi$（加法定数（ここでは+50）を除いて定義される）である。視覚的な目的のため、両方の Gauss 密度は100倍されている。}
\end{figure}

\begin{figure}[h!]
\centering
\begin{tabular}{@{}c@{\hspace{1mm}}c@{}}
\includegraphics[width=.35\linewidth]{gaussian-1d/interp-density}
\includegraphics[width=.25\linewidth]{gaussian-1d/gaussian-param}
\end{tabular}
\caption{\label{fig-1d-gaussian}
2つの1次元 Gauss 分布間の変位補間の計算。
%
Gauss 密度を $\Gg_{m,\si}(x) \eqdef \frac{1}{\sqrt{2\pi}s}e^{-\frac{(x-m)^2}{2s^2}}$ と表記して、補間 $\Gg_{(1-t)m_0+t m_1,(1-t)\si_0+t \si_1}$ を示す。
}
\end{figure}
